\documentclass[a4paper]{article}     
\usepackage{amsfonts}
\usepackage{amsmath}
\usepackage{amssymb}
\usepackage{graphicx}

\newcommand{\Bgcot}{\operatorname{Bgcot}}

\begin{document}

\noindent \large Auteur: Mark Gijsbrechts \\
\noindent \large Studierichting: Informatica\\
\noindent \large Oefeningenreeks 2B nr. 3.4\\

\medskip

\normalsize



Toon aan: $\Bgcot x + \Bgcot (-x) = \pi$ \\

$\alpha = \Bgcot x$ met $\alpha \in \left] 0, \pi \right[$ \\

$\Leftrightarrow x = \cot \alpha$ \\

$ -x = -\cot \alpha$ \\

$ -x = \cot(\pi - \alpha)$ met $\pi - \alpha \in \left] 0, \pi \right[$ \\

$\Leftrightarrow \pi - \alpha = \Bgcot(-x)$ \\ \\

$\Bgcot x + \Bgcot (-x) = \alpha + \left(\pi - \alpha\right) = \pi$ \\


\begin{thebibliography}{999}
%\bibitem{a} Referentie
\end{thebibliography}
\end{document} 
